% ===========================================================================
% circularity.tex -- the residual definition's loop, and the ordering that cuts
% it. \input by Sections/04c-reencoding.tex at sec:methods-residual. Compiles
% against thesis_v2/preamble.tex and nothing else: no external data, no
% \includegraphics, no package this document does not already load
% (tikz + arrows.meta, both in preamble.tex section 9).
%
% THIS IS A PROVENANCE DIAGRAM, NOT A RESULT. Deliberately absent: every NIR
% value, every arm, every split-level number. Two counts appear, once, and both
% are already in the prose of this same section. Sibling circularity.json lists
% every number drawn.
%
% THE LEFT PANEL CARRIES NO BUILD LABEL, DELIBERATELY. It drew "THE EARLIER
% BUILD -- NOT THE SHIPPED ONE" until this revision, which spent half a
% full-width figure on a pipeline the thesis did not run. The historical fact
% -- that an earlier build did close this loop -- is a disclosure, and a
% disclosure belongs in prose where it can be qualified: it is stated at
% Sections/04c-reencoding.tex:148-151 ("That is a loop [...], and an earlier
% build closed it, with generic and filter estimated over every condition and
% the holdout assigned afterwards"). That sentence stands on its own and was
% NOT deleted when the label came off the panel. What the left panel draws now
% is the structural hazard itself, in the present tense and attached to no
% build: the cycle the four definitions form, and the edge that fitting before
% splitting adds to it. The two panels therefore read hazard and remedy, which
% is the only reason the staged ordering is load-bearing rather than fastidious.
%
%   2,523 of 4,999 training conditions retained (50.5%)   04c-reencoding.tex:211-212
%   the resolved reliability line, +0.1086                 04c-reencoding.tex:211
%
% Cross-checked, all four agreeing, before drawing:
%   Sections/03-Apparatus.tex:206     "2,523 of 4,999 train (50.5\%)"
%   Sections/04c-reencoding.tex:212   "$2{,}523$ of $\num{4999}$ ... ($50.5\%$)"
%   Sections/04d-baselines.tex:337    "only $\num{2523}$ are training"
%   Sections/05-Limitations.tex:33    "$\num{2523}$ of $\num{4999}$ ... $\cos > 0.109$"
% +0.109 is the measured null's 95th percentile and +0.1086 is what it resolves
% to; the prose prints both, and the node carries the resolved value because
% that is the one the builder applies.
%
% LAYOUT NOTES, all forced by a build:
%   * fullwidth = 174mm = 493.23pt (textwidth 142mm + marginparsep 4mm +
%     marginparwidth 28mm, preamble.tex:88-90 and :525). The bounding box is
%     PINNED to 17.35cm = 491.81pt, 1.4pt inside the measure. Unpinned, the
%     picture measured 3.10pt OVER and the figure body reported an overfull
%     hbox: nodes with anchor=north west hang their 3.33pt inner sep left of
%     the anchor at x=0.20, and the strip rules end at x=17.35. Pinning clips
%     no ink -- no drawn object lies outside [0, 17.35] x [-0.95, 7.95].
%   * NO node uses text width. Every line break is written by hand. The first
%     build did use text width and TeX hyphenated inside the boxes
%     ("the relia-/bility line"), grew four of them past their minimum heights
%     and overlapped the rows below. Hand-broken lines make every box width and
%     height deterministic, which is what lets the coordinates below be trusted.
%   * the four objects are ONE node set drawn twice, same style, same wording,
%     same 2.80cm minimum width, so the right panel reads as a re-laying of the
%     left and not as a second diagram. Only the arrangement changes.
%   * the fatal edge is drawn and STRUCK, not deleted. A deleted edge leaves
%     nothing for the caption to point at, and the passage is about which single
%     edge the ordering removes.
%   * the two streams leaving SPLIT must not cross. They cannot both turn upward
%     in the order they leave: the stream that leaves HIGHER (held-out, 2.2mm
%     above the east anchor) turns up FIRST (x=10.30) and the stream that leaves
%     LOWER (training keys, at the east anchor) turns up LATER (x=10.85). Any
%     other pairing crosses, and a crossing here reads as contact between them.
%   * both streams use -| rather than an explicit corner coordinate, so the
%     corner inherits the node anchor's own y. Written as a literal corner, the
%     first segment came out as a visible diagonal whenever a node's height
%     differed by a hair from the value assumed at design time.
%   * the held-out stream is routed over the top of the fit column, and its one
%     arrowhead points right, into TRANSFORM. No arrowhead anywhere in the right
%     panel points back left or up into the fit; that absence is the claim.
%   * NO edge label is masked onto its arrow. The first build masked all of
%     them and the labels ate the edges: the C->D gap in the left panel is
%     1.20cm and "subtracted to give" is 1.15cm wide, so the bottom edge of the
%     cycle disappeared and the loop stopped reading as a loop. Every label now
%     sits beside its edge. The tightest is "training keys", hung under the
%     corner at x = 10.43 to 11.27, between the held-out riser at 10.30 and the
%     fit column's left edge at 11.35.
% ===========================================================================
\begin{figure}[htbp]
\begin{fullwidth}
\centering
\begin{adjustbox}{max width=\linewidth}
\begin{tikzpicture}[x=1cm, y=1cm]
  \useasboundingbox (0,-0.95) rectangle (17.35,7.95);
  \tikzset{
    % the shared node set. Identical style in both panels.
    obj/.style={draw=rulegrey, line width=0.4pt, fill=white, inner sep=3pt,
                align=center, font=\scriptsize, text=ink,
                minimum width=2.80cm, minimum height=0.68cm},
    edg/.style={-{Stealth[length=3.6pt,width=2.8pt]}, draw=ink,
                line width=0.5pt},
    lab/.style={font=\scriptsize, text=quietink, fill=white, inner sep=1pt,
                align=center},
    side/.style={font=\scriptsize, text=quietink, align=flush right,
                 anchor=east, text width=3.20cm},
    quiet/.style={font=\scriptsize\itshape, text=quietink, align=center},
    tag/.style={font=\sffamily\scriptsize\bfseries, text=accent},
    sub/.style={font=\footnotesize\itshape, text=quietink}}
  \def\T#1{{\sffamily\scriptsize\bfseries\color{accent}#1}\\[1.0pt]}
  \def\Q#1{\\[1.0pt]{\color{quietink}#1}}

  % ======================= PANEL HEADS =====================================
  % Hazard and remedy, one head each, same shape: a tag naming the structure
  % and a sub saying what happens to the four objects. Neither names a build.
  \node[tag, anchor=north west] at (0.20,7.90) {THE CIRCULARITY};
  \node[sub, anchor=north west] at (0.20,7.55)
    {the four objects define one another};
  \node[tag, anchor=north west] at (7.50,7.90) {THE ORDERING THAT CUTS IT};
  \node[sub, anchor=north west] at (7.50,7.55)
    {the same four objects, staged by execution order};

  \draw[rulegrey, line width=0.4pt] (7.20,1.15) -- (7.20,6.60);

  % ======================= LEFT: THE CLOSED LOOP ===========================
  \node[obj, anchor=north] (A) at (1.60,6.10) {the reliability line};
  \node[obj, anchor=north] (B) at (5.60,6.10) {the training set};
  \node[obj, anchor=north] (C) at (5.60,4.45) {the generic};
  \node[obj, anchor=north] (D) at (1.60,4.45) {the residual};

  % NONE of the four loop labels is masked onto its arrow. Masked, they ate the
  % edges: the C->D gap is 1.20cm and "subtracted to give" is 1.15cm wide, so
  % the bottom edge of the cycle vanished entirely and the loop stopped reading
  % as a loop. Each label now sits beside its own edge -- the two horizontals
  % labelled outside the cycle, the two verticals inside it -- and all four
  % edges draw unbroken.
  \draw[edg] (A.east) -- (B.west);
  \draw[edg] (B.south) -- (C.north);
  \draw[edg] (C.west) -- (D.east);
  \draw[edg] (D.north) -- (A.south);
  \node[lab, anchor=south] at (3.60,5.86) {selects};
  \node[lab, anchor=east]  at (5.45,4.94) {fits};
  % clears the box row's bottom edge at y=3.77: at y=3.97 its first line sat
  % level with the two boxes' lower corners, 0.02cm from each, and read as a
  % collision even though it was not one.
  \node[lab, anchor=north] at (3.60,3.88) {subtracted\\to give};
  \node[lab, anchor=west]  at (1.75,4.94) {must reproduce\\across a half-split};

  % --- the fatal edge: drawn once, struck once, and labelled --------------
  % Present tense, and attached to an ordering rather than to a build: this is
  % what fitting before splitting does to the cycle, not what one run once did.
  % The quiet italic under H answers "why would they enter the fit?" inside the
  % panel, and mirrors the right panel's one quiet italic under SPLIT, so each
  % half carries exactly one such annotation.
  \node[obj, anchor=north] (H) at (1.60,2.60) {held-out conditions};
  \draw[edg] (H.east) -| (C.south);
  \node[lab, anchor=east] at (5.45,2.95) {enter the fit};
  % Its y is the strike caption's y, not a value of its own: both are two-line
  % blocks of the same size, so sharing anchor y = 1.74 gives them a common top
  % AND a common bottom (0.944) and the bottom-left reads as one row rather than
  % two blocks at two heights. Widest line measures 54.42pt = 1.91cm, so with
  % inner sep the node spans x 0.53-2.67, clear of the strike caption's left
  % edge at x = 3.10 by 0.43cm. Top sits 0.18cm under H, matching the 0.20cm
  % the right panel leaves under SPLIT.
  \node[quiet, anchor=north] at (1.60,1.74)
    {when the fit\\precedes the split};
  % the strike. One stroke. Nothing else in this figure is accent-coloured
  % except the structural tags.
  \draw[accent, line width=1.1pt] (4.25,1.94) -- (4.75,2.58);
  \node[font=\sffamily\scriptsize, text=accent, anchor=north west,
        align=left] at (3.10,1.74)
    {struck: the one edge\\the ordering removes};

  % ======================= RIGHT: THE STAGED PIPELINE ======================
  \fill[panelbg] (7.45,3.55) rectangle (10.15,6.30);
  \node[obj, anchor=north, minimum width=2.55cm] (INV) at (8.80,6.10)
    {\T{INVENTORY}wells, drugs, doses};
  \node[obj, anchor=north, minimum width=2.55cm] (SPL) at (8.80,5.10)
    {\T{SPLIT}train \textbar{} held out};
  \draw[edg] (INV.south) -- (SPL.north);
  \node[quiet, anchor=north] at (8.80,4.22)
    {metadata only:\\no expression data};

  % FIT, and the training side that descends from it
  \node[obj, anchor=north] (G) at (12.75,6.10) {\T{FIT}the generic};
  \node[obj, anchor=north] (R) at (12.75,4.60) {the residual};
  \node[obj, anchor=north] (L) at (12.75,3.30)
    {the reliability line\Q{$+0.1086$}};
  \node[obj, anchor=north] (S) at (12.75,1.90) {the training set};
  \draw[edg] (G.south) -- (R.north);
  \draw[edg] (R.south) -- (L.north);
  \draw[edg] (L.south) -- (S.north);

  \node[side] at (11.25,2.80) {$2{,}523$ of $4{,}999$ training conditions
    retained ($50.5\%$)};
  \node[side] at (11.25,1.56) {the residual targets the model is trained on};

  % training keys, handed to the fit explicitly
  % -| , not an explicit corner coordinate: the corner must inherit SPL.east's
  % own y, or the first segment comes out as a visible diagonal when the node's
  % height differs by a hair from the value assumed here.
  % This is the one edge label NOT masked onto its arrow. The riser is 1.0cm
  % long and a masked label ate 0.6cm of it, leaving the arrow looking like two
  % disconnected stubs. It hangs under the corner instead, in the only clear
  % space in the panel: x 10.43-11.27, below the horizontal and above the
  % retention annotation, which tops out at y = 3.25.
  \draw[edg] (SPL.east) -| (10.85,5.675) -- (G.west);
  \node[font=\scriptsize, text=quietink, align=center, anchor=north]
    at (10.85,4.50) {training\\keys};

  % TRANSFORM, fed by the fitted generic and by a stream that bypasses the fit
  \node[obj, anchor=north] (TR) at (15.95,6.10)
    {\T{TRANSFORM}held-out conditions\\projected through\\the fitted generic};
  \draw[edg] (G.east) -- (G.east -| TR.west);
  \draw[edg] ([yshift=2.2mm]SPL.east) -| (10.30,6.45) -- (15.95,6.45)
    -- (TR.north);
  \node[font=\scriptsize, text=quietink, anchor=south] at (13.10,6.55)
    {held-out keys --- they never enter the fit};

  \node[obj, anchor=north] (HT) at (15.95,4.30)
    {held-out residual\\targets};
  \draw[edg] (TR.south) -- (HT.north);
  \node[quiet, anchor=north] at (15.95,3.25)
    {deliberately unfiltered};
  \node[quiet, anchor=north] at (15.95,2.80)
    {nothing here\\re-enters the fit};

  % ======================= THE POISON TEST STRIP ===========================
  \draw[rulegrey, line width=0.4pt] (0.20,0.85) -- (17.35,0.85);
  \node[tag, anchor=north west, align=left] at (0.20,0.46)
    {THE POISON\\TEST};
  \node[font=\sffamily\scriptsize\bfseries, text=ink, anchor=north west]
    at (3.30,0.55) {CORRUPTION APPLIED};
  \node[font=\sffamily\scriptsize\bfseries, text=ink, anchor=north west]
    at (9.60,0.55) {REQUIRED OUTCOME};
  \draw[rulegrey, line width=0.4pt] (3.30,0.20) -- (17.35,0.20);
  \node[font=\sffamily\footnotesize, text=ink, anchor=north west]
    at (3.30,0.08) {corrupt every held-out expression vector};
  \node[font=\sffamily\footnotesize, text=ink, anchor=north west]
    at (9.60,0.08) {the training targets are byte-identical};
  \node[font=\sffamily\footnotesize, text=ink, anchor=north west]
    at (3.30,-0.40) {corrupt one training condition};
  \node[font=\sffamily\footnotesize, text=ink, anchor=north west]
    at (9.60,-0.40) {the training targets must move};
  \draw[rulegrey, line width=0.4pt] (3.30,-0.85) -- (17.35,-0.85);
\end{tikzpicture}
\end{adjustbox}
\caption[The residual's circular definition, cut by
ordering]{\captiontitle{The definition of the residual could have closed on itself,
and what stops it is execution order, not an argument.} \emph{Left, the
circularity.} A condition enters training only if
its residual reproduces across a half-split; the residual exists only once the
generic is subtracted; the generic is fitted on the training set; and the
training set is what the reliability line selects. Estimate the generic and the
filter over every condition and assign the holdout afterwards, and held-out
conditions enter the fit --- the struck edge --- so each held-out
target is defined partly by other held-out conditions. \emph{Right, the
ordering that cuts it,} the build this thesis reports. The same four objects,
staged. Inventory and split read
metadata only, before any pseudobulk exists, so nothing downstream of the split
can have defined it; the fit is handed the training keys explicitly; and
transform projects held-out conditions through the fitted generic in one
direction, with no path back into it. The generic is the mean over training
drugs other than this one, within plate. The two counts are the
\emph{training-set} retention at the resolved reliability line $+0.1086$, the
$95$th percentile of a null pairing half $A$ of one condition with half $B$ of
another; the held-out sets are deliberately unfiltered, since dropping held-out
conditions for failing their own reproducibility bar would select the
evaluation set on its own outcome. \emph{Beneath,} the two assertions that make
the ordering checkable in both directions rather than merely asserted. This is
a provenance diagram: no $\NIR$ value, arm or split-level result appears in
it.}
\label{fig:circularity}
\end{fullwidth}
\end{figure}
